% GENERATED FILE. Edit canonical lessons, then run npm run generate:books.
% canonical-source-hash: fnv1a64:af1b389838d5ebb2
% canonical-lessons: ZH-C141-luyou, ZH-C141-youxi, ZH-C141-guangjie, ZH-C141-sanbu, ZH-C141-diaoyu

\chapter{Travel and Fun}
\label{ch:zh-travel-and-fun}
\hlchaptermodality{141}

\begin{chapteropening}
\textbf{By the end of this chapter:} \emph{I can talk about travelling, games, going round the shops, walks and fishing.}

Complete the last lesson of chapter 141: I can talk about travelling, games, going round the shops, walks and fishing.
\end{chapteropening}

\section[lǚyóu]{lǚyóu — to travel; to tour}

\label{lesson:ZH-C141-luyou}



\begin{quote}
Before the new one: say the Mandarin for most, then the Mandarin for to calculate.
\end{quote}

\subsection*{You'll want to know: lǚyóu}
\emph{lǚyóu} — \textquotedblleft{}to travel; to tour\textquotedblright{}.

\begin{quote}
\emph{wǒ xǐhuan lǚyóu} — I love travelling.
\end{quote}

Its characters come later; for now it is a word for the ear.

\begin{cousinweb}
\emph{lǚyóu} is built from \emph{lǚ} and \emph{yóu}, both inherited from Middle Chinese, \emph{ljoX} and \emph{yuw} (Guangyun, in Baxter's transcription). The dictionary glosses the first 'a body of troops' and the second 'to float, to roam', the \emph{yóu} of \emph{yóuyǒng}. 'Journey' for the first is not given there; how it got that sense is not settled.
\end{cousinweb}

\begin{grammarlens}[title={What you've built}]
One more word for what I like.
\end{grammarlens}

\subsection*{Your turn}
\begin{itemize}
\raggedright
  \item \emph{Say it:} \emph{lǚyóu}
  \item \emph{Say it:} \emph{lǚyóu}, once more
  \item \emph{Say it:} \emph{wǒ xǐhuan lǚyóu} — I love travelling.
  \item \emph{Recall:} say the Mandarin for most, then the Mandarin for to calculate, then say \emph{lǚyóu} again
\end{itemize}

\begin{tcolorbox}[breakable,colback=teal!4,colframe=teal!35!black,title={Before you move on}]
What does lǚyóu mean? (\textquotedblleft{}To travel; to tour\textquotedblright{}.) Say it once more.
\end{tcolorbox}
\section[yóuxì]{yóuxì — game}

\label{lesson:ZH-C141-youxi}



\begin{quote}
Before the new one: say the Mandarin for to calculate, then the Mandarin for to travel.
\end{quote}

\subsection*{You'll want to know: yóuxì}
\emph{yóuxì} — \textquotedblleft{}game\textquotedblright{}.

\begin{quote}
\emph{wǒ érzi hěn xǐhuan wán yóuxì} — My son loves playing games.
\end{quote}

Its characters come later; for now it is a word for the ear.

\begin{cousinweb}
\emph{yóuxì} is built from \emph{yóu} 'to roam' and \emph{xì} 'to play, to tease', both inherited from Middle Chinese, \emph{yuw} and \emph{xjeH} (Guangyun, in Baxter's transcription).
\end{cousinweb}

\begin{grammarlens}[title={What you've built}]
One more word for what I like.
\end{grammarlens}

\subsection*{Your turn}
\begin{itemize}
\raggedright
  \item \emph{Say it:} \emph{yóuxì}
  \item \emph{Say it:} \emph{yóuxì}, once more
  \item \emph{Say it:} \emph{wǒ érzi hěn xǐhuan wán yóuxì} — My son loves playing games.
  \item \emph{Recall:} say the Mandarin for to calculate, then the Mandarin for to travel, then say \emph{yóuxì} again
\end{itemize}

\begin{tcolorbox}[breakable,colback=teal!4,colframe=teal!35!black,title={Before you move on}]
What does yóuxì mean? (\textquotedblleft{}Game\textquotedblright{}.) Say it once more.
\end{tcolorbox}
\section[guàngjiē]{guàngjiē — to stroll the shops; to go window-shopping}

\label{lesson:ZH-C141-guangjie}



\begin{quote}
Before the new one: say the Mandarin for to travel, then the Mandarin for a game.
\end{quote}

\subsection*{You'll want to know: guàngjiē}
\emph{guàngjiē} — \textquotedblleft{}to stroll the shops; to go window-shopping\textquotedblright{}.

\begin{quote}
\emph{zhōumò wǒmen qù guàngjiē} — At the weekend we're going round the shops.
\end{quote}

Its characters come later; for now it is a word for the ear.

\begin{cousinweb}
\emph{guàngjiē} is built from \emph{guàng} and \emph{jiē} 'street', both inherited from Middle Chinese, \emph{kjwangX} and \emph{kea} (Guangyun, in Baxter's transcription). The dictionary glosses the first 'the look of walking'; the regular sound change would give it a third tone, so the fourth tone of \emph{guàng} is not settled.
\end{cousinweb}

\begin{grammarlens}[title={What you've built}]
One more word for what I like.
\end{grammarlens}

\subsection*{Your turn}
\begin{itemize}
\raggedright
  \item \emph{Say it:} \emph{guàngjiē}
  \item \emph{Say it:} \emph{guàngjiē}, once more
  \item \emph{Say it:} \emph{zhōumò wǒmen qù guàngjiē} — At the weekend we're going round the shops.
  \item \emph{Recall:} say the Mandarin for to travel, then the Mandarin for a game, then say \emph{guàngjiē} again
\end{itemize}

\begin{tcolorbox}[breakable,colback=teal!4,colframe=teal!35!black,title={Before you move on}]
What does guàngjiē mean? (\textquotedblleft{}To stroll the shops; to go window-shopping\textquotedblright{}.) Say it once more.
\end{tcolorbox}
\section[sànbù]{sànbù — to take a walk}

\label{lesson:ZH-C141-sanbu}



\begin{quote}
Before the new one: say the Mandarin for a game, then the Mandarin for to stroll the shops.
\end{quote}

\subsection*{You'll want to know: sànbù}
\emph{sànbù} — \textquotedblleft{}to take a walk\textquotedblright{}.

\begin{quote}
\emph{wǎnfàn yǐhòu wǒmen qù sànbù} — After dinner we go for a walk.
\end{quote}

Its characters come later; for now it is a word for the ear.

\begin{cousinweb}
\emph{sànbù} is built from \emph{sàn} 'to scatter, to spread out' and \emph{bù} 'step', both inherited from Middle Chinese, \emph{sanH} and \emph{buH} (Guangyun, in Baxter's transcription): steps let loose.
\end{cousinweb}

\begin{grammarlens}[title={What you've built}]
One more word for what I like.
\end{grammarlens}

\subsection*{Your turn}
\begin{itemize}
\raggedright
  \item \emph{Say it:} \emph{sànbù}
  \item \emph{Say it:} \emph{sànbù}, once more
  \item \emph{Say it:} \emph{wǎnfàn yǐhòu wǒmen qù sànbù} — After dinner we go for a walk.
  \item \emph{Recall:} say the Mandarin for a game, then the Mandarin for to stroll the shops, then say \emph{sànbù} again
\end{itemize}

\begin{tcolorbox}[breakable,colback=teal!4,colframe=teal!35!black,title={Before you move on}]
What does sànbù mean? (\textquotedblleft{}To take a walk\textquotedblright{}.) Say it once more.
\end{tcolorbox}
\section[diàoyú]{diàoyú — to go fishing}

\label{lesson:ZH-C141-diaoyu}



\begin{quote}
Before the new one: say the Mandarin for to stroll the shops, then the Mandarin for to take a walk.
\end{quote}

\subsection*{You'll want to know: diàoyú}
\emph{diàoyú} — \textquotedblleft{}to go fishing\textquotedblright{}.

\begin{quote}
\emph{wǒ bàba xǐhuan diàoyú} — My dad loves fishing.
\end{quote}

Its characters come later; for now it is a word for the ear.

\begin{cousinweb}
\emph{diàoyú} is built from \emph{diào} 'to angle' and \emph{yú} 'fish', both inherited from Middle Chinese, \emph{tewH} and \emph{ngjo} (Guangyun, in Baxter's transcription); the dictionary glosses the first with this very word, 'to fish'.
\end{cousinweb}

\begin{grammarlens}[title={What you've built}]
One more word for what I like.
\end{grammarlens}

\subsection*{Your turn}
\begin{itemize}
\raggedright
  \item \emph{Say it:} \emph{diàoyú}
  \item \emph{Say it:} \emph{diàoyú}, once more
  \item \emph{Say it:} \emph{wǒ bàba xǐhuan diàoyú} — My dad loves fishing.
  \item \emph{Recall:} say the Mandarin for to stroll the shops, then the Mandarin for to take a walk, then say \emph{diàoyú} again
\end{itemize}

\begin{tcolorbox}[breakable,colback=teal!4,colframe=teal!35!black,title={Before you move on}]
What does diàoyú mean? (\textquotedblleft{}To go fishing\textquotedblright{}.) Say it once more.
\end{tcolorbox}
